\section{Introduction}
\label{sec:intro}

Large-scale benchmark suites such as \helm \citep{liang2023helm} have become the
shared substrate on which the community compares language models. A central
promise of these suites is transparency: for many (scenario, model) pairs the
public corpus ships the exact prompts sent to the model, the completions it
returned, and the metrics computed from them. This apparent completeness invites
a natural reproducibility question --- \emph{re-run the evaluation locally and
check that the numbers match} --- and a natural expectation that open-weight
models, whose weights anyone can download, should reproduce cleanly.

Our starting point was exactly this program: reproduce a large slice of public
\helm on local open-weight hardware. It did not survive contact with the data,
and \emph{the way it failed is the contribution of this paper.} Two facts
dominate. First, a large fraction of the open-weight corpus was originally served
not by the harness itself but through a hosted inference provider (frequently
Together AI), whose software and hardware stack the public record does not
describe. Second, and more fundamentally, the artifact that \emph{does} survive
--- the resolved \file{run\_spec.json} --- captures \emph{what} to generate
(scenario, prompts, decoding parameters, metrics) and the \emph{name} of the
model and its deployment, but almost nothing about the \emph{execution substrate}:
the precision the weights were loaded at, the exact checkpoint revision, the
attention kernel, the tokenizer and chat-template versions, and the software-stack
versions. A benchmark recipe, in other words, is not a complete experimental
record.

Once this is taken seriously, ``did it reproduce?'' fractures into three distinct
questions that the literature often conflates: (i) \emph{can the evaluation be run
at all?} (gated models and datasets, revoked download endpoints, packaging
incompatibilities, unavailable judges --- these are eligibility and environment
filters, not reproduction outcomes); (ii) \emph{can a faithful local execution be
reconstructed} when the key execution parameters were never written down and the
original may have been served by an external provider; and (iii) when local and
public numbers differ, \emph{to what is the difference attributable} --- a genuine
recipe disagreement, a controllable execution-substrate parameter, an
artifact-processing effect, or an irreducibly lost provenance?

\paragraph{Thesis.} We therefore reframe benchmark reproduction as a layered
\emph{system-identification} problem. The measured score is a function
$Y = F(R, D, I, A, U)$ of the recorded recipe $R$, the deployment $D$, the
execution instrument $I$, the artifact-interpretation history $A$, and residual
stochastic/numerical influences $U$ (\S\ref{sec:model}). The public record
exposes only a projection of these variables. Controlled reconstruction can
\emph{attribute} an observed disagreement to a specific layer, and can sometimes
\emph{close} it by supplying the missing parameter; but when provenance has been
lost or transformed, the original measurement can be \emph{non-identifiable from
the surviving evidence examined}. This is a more defensible and more scientifically
interesting claim than a blanket ``open-weight \helm is (ir)reproducible.''

\paragraph{Contributions.} Concretely, this paper contributes:
\begin{enumerate}[nosep,leftmargin=1.4em]
  \item \textbf{A measurement model and a disagreement taxonomy}
        (\S\ref{sec:model}). We factor a benchmark result into recipe,
        deployment, execution-instrument, artifact-interpretation, and residual
        layers; classify observed disagreements into six causes; and separate
        that ``why'' taxonomy from three orthogonal reproduction \emph{targets}
        --- artifact reconstruction, procedural reproduction, and claim-level
        robustness --- clarifying what ``reproduction'' should even mean.
  \item \textbf{\textsc{EvalAudit}, a reproducibility-auditing system}
        (\S\ref{sec:system}): exact \file{run\_spec.json} replay (never recipe
        reconstruction), an execution-substrate reconstruction and search tool,
        version-pinned ``era'' scoring containers that freeze the measurement
        instrument, and a framework-free layered comparison core.
  \item \textbf{Attribution case studies on modern open-weight models}
        (\S\ref{sec:olmo}). On OLMo-family models we attribute concrete
        disagreements to a tokenizer special-token append, to a chat-template
        \code{add\_generation\_prompt} flag whose behavior depends on the
        \code{transformers} version, and to load precision --- while being
        explicit that the precision result is at present a configuration
        \emph{discovery} awaiting ordinary-path confirmation, and that one dense
        OLMo-2 divergence remains unresolved.
  \item \textbf{A recoverable-vs-non-identifiable contrast}
        (\S\ref{sec:redpajama}--\ref{sec:classic}). RedPajama-3B demonstrates
        end-to-end execution under declared historical proxy instruments, whereas
        the original-paper classic models (GPT-J, GPT-NeoX, OPT) carry a metric
        class path that resolves in \emph{no} released \helm version and
        triangulates to unreleased pre-v0.1.0 code, making the producing
        instrument non-identifiable from surviving evidence.
  \item \textbf{A minimal provenance standard} (\S\ref{sec:provenance}). We show
        that three recorded fields --- \code{model\_revision},
        \code{torch\_dtype}, and \code{transformers\_version} --- transitively
        close most of the identifiability gap, and give a backward-compatible way
        to record them (outside the run-spec identity string, so official$
        \leftrightarrow$local pairing is preserved), plus an artifact-release
        checklist.
\end{enumerate}

\paragraph{Scope and honesty about evidence.} This paper is deliberately explicit
about the maturity of each empirical claim. Several results rest on
configuration-discovery probes rather than completed end-to-end reproductions;
we mark these throughout and describe the confirmation experiment that would
promote them (\S\ref{sec:olmo-confirm}). The system, the conceptual framework,
and the recoverable-vs-non-identifiable contrast stand independently of whether
any single confirmation succeeds.

This work builds directly on \eee \citep{eee2026}, a normalized artifact
representation that \emph{detects} reproducibility discrepancies but, by its
authors' own statement, cannot separate checkpoint, quantization, precision,
tokenizer, and generation-kernel effects. Our positioning is complementary and
one sentence long: \emph{\eee detects reproducibility gaps; this work studies
their attribution, possible closure, and identifiability.}
